% !TEX root = ../proyect.tex

\chapter*{Resumen}\label{cap:resumen}
\addcontentsline{toc}{chapter}{Resumen}

El acceso a la información académica en una universidad presencial sigue siendo costoso para el alumnado. Los datos del día a día (asignaturas, horarios, profesorado, evaluación) se reparten entre planes docentes en PDF, portales como Sevius y páginas de departamentos, y el alumno no siempre sabe en qué portal buscar cada cosa. El problema afecta especialmente a los de nuevo ingreso e internacionales, pero también al alumnado matriculado: comparar criterios de evaluación obliga a abrir un PDF por asignatura, contactar con un profesor exige cruzar varios directorios y la consulta de horarios resulta confusa cuando grupo y cuatrimestre no se identifican con claridad.

Este Trabajo Fin de Grado presenta \textbf{Linceus Assistant}, un asistente conversacional para el alumnado de la ETSII de la Universidad de Sevilla que centraliza esa información en un único punto. Frente a los chatbots universitarios desplegados en el ecosistema español, la propuesta se diferencia en dos aspectos. Primero, está enfocada al \textbf{estudiante ya matriculado}, no al ciclo de captación y matrícula que cubren los asistentes institucionales existentes. Segundo, adopta una arquitectura de \textbf{recuperación aumentada por generación} (RAG) sobre los documentos oficiales, lo que le permite responder a preguntas no anticipadas en un guion cerrado, frente al modelo de menús predefinidos habitual en las soluciones comerciales.

El sistema se ha construido sobre Rasa~3.6, integrado con la API de Gemini para clasificación de intenciones, generación de SQL y renderizado de respuestas. La base de datos reside en Supabase con la extensión \texttt{pgvector}, lo que permite combinar consultas estructuradas con recuperación semántica sobre los planes docentes. El despliegue se realiza con Docker Compose en cuatro contenedores. El producto cubre los tres dominios funcionales del prototipo (\textbf{asignaturas, horarios y profesorado}), incorpora cambio dinámico de titulación y dispone de un panel de administración para garantizar la mantenibilidad del sistema.

La validación se ha organizado en cinco capas independientes y todas superan su umbral académico de referencia: \textbf{94,3\,\%} en la suite E2E (159 casos), \textbf{91,9\,\%} en RAG sobre 74 preguntas extraídas de planes docentes reales, \textbf{100\,\%} en política Rasa, y \textbf{98,3\,\%} en el piloto con 59 sesiones de alumnos reales. Las métricas no funcionales acompañan: latencia mediana de 3,4~segundos por turno y coste extrapolado entre 8~y 50~euros anuales para 700~alumnos, dos órdenes de magnitud por debajo de cualquier suscripción comercial equivalente.

La arquitectura es \textbf{extensible y escalable} por construcción: ampliar el sistema a nuevas titulaciones u otros centros de la Universidad de Sevilla, o incorporar dominios adicionales, descansa principalmente en la carga de datos a través del panel de administración, sin tocar el núcleo conversacional. El proyecto demuestra que es viable construir un asistente universitario abierto, validado cuantitativamente y económicamente sostenible empleando herramientas accesibles a un trabajo de fin de grado.
